\input{preamble}

\begin{document}

\title{Partial Differential Equations}
\maketitle

\phantomsection
\label{section-phantom}

\tableofcontents

\section{Introduction}
\label{section-introduction}

\noindent
These notes develop the theory of partial
differential equations.

Given the material elsewhere in the Heap
Project, the basic things to learn are:

\begin{enumerate}
\item The model elliptic, parabolic, and
hyperbolic equations: the Laplace, heat,
and wave equations.

\item How characteristic curves explain
first-order equations, transport, and the
formation of singularities.

\item The fundamental solutions and Green
functions of the model equations, together
with the maximum principle and mean-value
property.

\item Fourier-series and Fourier-transform
methods, separation of variables, and the
role of distributions.

\item Sobolev spaces, weak derivatives,
embeddings, compactness, and trace
theorems.

\item Weak formulations and the main
existence mechanisms: energy estimates,
the Lax--Milgram theorem, compactness, and
the calculus of variations.

\item Elliptic regularity: when a weak
solution becomes smooth, and how the
regularity of the data controls that of
the solution.

\item Energy and semigroup methods for
evolution equations, including existence,
uniqueness, stability, and long-time
behavior.

\item The distinction between classical,
weak, distributional, and viscosity
solutions, and why different equations
require different notions of solution.

\item How differential operators fit the
geometric language already developed in
the project: symbols, ellipticity, the
Laplace--Beltrami operator, and de Rham
and Hodge Laplacians.

\item The first nonlinear phenomena:
conservation laws, shocks, blow-up,
scaling, and critical exponents.
\end{enumerate}

\section{Some Problems in Hamiltonian
Geometry of PDEs}
\label{section-hamiltonian-geometry-pdes}

\medskip\noindent
{\bf Speaker.}

Raffaele Vitolo,
Universit\`a del Salento.

\medskip\noindent
{\bf Seminar information.}

Semin\'ario das Sextas, PUC-Rio,
September 27, 2024. The talk was also
presented in the Geometric Structures on
Manifolds Seminar at IMPA on
December 5, 2024.

\medskip\noindent
{\bf Abstract.}

The Hamiltonian formulation of partial
differential equations is one of the
cornerstones of the theory of integrable
systems. It is modeled on finite-dimensional
Hamiltonian systems and has an intrinsic
geometric nature. The talk reviews
differential-geometric aspects of the
Hamiltonian theory of PDEs and discusses
projective-geometric properties of
integrable systems that have emerged in
recent years.

\medskip\noindent
{\bf Hamiltonian evolutionary equations.}

Consider an evolutionary system
$$
u_t^i=f^i(t,x,u^j,u_x^j,u_{xx}^j,\ldots).
$$
A local functional has the form
$$
\mathcal{H}=\int h\,dx.
$$
Its variational derivative is
$$
\frac{\delta\mathcal{H}}{\delta u^i}
=
\sum_\sigma(-D)_\sigma
\frac{\partial h}{\partial u^i_\sigma}.
$$

\begin{definition}
\label{definition-hamiltonian-pde}
The system has a \textit{Hamiltonian
formulation} if there is a differential
operator $A$ and a local functional
$\mathcal{H}$ such that
$$
u_t^i=A^{ij}
\frac{\delta\mathcal{H}}{\delta u^j}.
$$
\end{definition}

The operator $A$ plays the role of a
Poisson tensor. It sends the differential
of a Hamiltonian functional to the
corresponding evolutionary vector field.

\begin{definition}
\label{definition-hamiltonian-operator}
A differential operator $A$ is a
\textit{Hamiltonian operator} if
$$
\{\mathcal{F},\mathcal{G}\}_A
=
\int
\frac{\delta\mathcal{F}}{\delta u^i}
A^{ij}
\frac{\delta\mathcal{G}}{\delta u^j}
\,dx
$$
defines a Poisson bracket on local
functionals.
\end{definition}

Thus $A$ must be formally skew-adjoint and
must satisfy the differential analogue of
the Jacobi identity.

There is also a purely algebraic definition
of differential operators. Let $R$ be a
ring, let $P,Q$ be $R$-modules, and let
$\Delta:P\to Q$ be additive. For $a\in R$,
write
$$
\delta_a\Delta=a\Delta-\Delta a.
$$
The operator $\Delta$ has order at most
$k$ if
$$
\delta_{a_0}\delta_{a_1}\cdots
\delta_{a_k}\Delta=0
$$
for all $a_0,\ldots,a_k\in R$. This is the
Grothendieck definition of a differential
operator.

\medskip\noindent
{\bf The KdV equation.}

\begin{example}
\label{example-kdv-bihamiltonian}
Consider the Korteweg--de Vries equation
$$
u_t=uu_x+u_{xxx}.
$$
It admits two Hamiltonian descriptions.
Set
$$
A_1=D_x
$$
and
$$
A_2=D_x^3+\frac{2}{3}uD_x
+\frac{1}{3}u_x.
$$
For the Hamiltonians
$$
\mathcal{H}_1=
\int\left(
\frac{u^3}{6}-\frac{u_x^2}{2}
\right)dx
$$
and
$$
\mathcal{H}_2=\int\frac{u^2}{2}\,dx,
$$
one has
$$
u_t=A_1
\frac{\delta\mathcal{H}_1}{\delta u}
=A_2
\frac{\delta\mathcal{H}_2}{\delta u}.
$$
\end{example}

This is the model for a bi-Hamiltonian
system.

\begin{definition}
\label{definition-compatible-hamiltonian}
Two Hamiltonian operators $A_1,A_2$ are
\textit{compatible} if every linear
combination of them is again Hamiltonian.
A system admitting such a pair is
\textit{bi-Hamiltonian}.
\end{definition}

\medskip\noindent
{\bf The Magri scheme.}

A Hamiltonian operator sends conservation
laws to symmetries. Given compatible
Hamiltonian operators, the Magri recursion
seeks functionals $\mathcal{H}_r$ satisfying
$$
A_1
\frac{\delta\mathcal{H}_{r+1}}{\delta u}
=
A_2
\frac{\delta\mathcal{H}_r}{\delta u}.
$$
Under suitable hypotheses, this produces
an infinite hierarchy of commuting flows
and conserved Hamiltonians in involution:
$$
\{\mathcal{H}_r,\mathcal{H}_s\}=0.
$$

This motivates one notion of integrability
for PDEs: the existence of infinitely many
commuting symmetries or conservation laws.
There is no direct infinite-dimensional
version of the Liouville theorem that
settles the theory, but the hierarchy may
lead to suitable coordinates or to an
inverse-scattering construction.

\medskip\noindent
{\bf Dubrovin--Novikov operators.}

A first-order homogeneous operator has the
form
$$
A^{ij}=g^{ij}(u)D_x
+b^{ij}_k(u)u_x^k.
$$
The conditions for $A$ to be Hamiltonian
have a differential-geometric meaning.
Formal skew-adjointness forces $g^{ij}$ to
be symmetric and relates $b^{ij}_k$ to a
metric connection. The Jacobi identity
then forces that connection to be
torsion-free and the metric to be flat,
when the metric is nondegenerate.

Thus the classification of these
Hamiltonian operators becomes a problem
in the geometry of flat metrics. The
classification of compatible pairs is
more rigid and encodes whole hierarchies
of conservation laws.

\medskip\noindent
{\bf Projective classification.}

The talk described a classification program
using projective reciprocal
transformations. Such transformations can
preserve the order and Hamiltonian nature
of an operator while changing its
coefficients. One first fixes a
higher-order operator and then classifies
compatible lower-order operators modulo
this group.

In the cases discussed in the talk,
second-order homogeneous Hamiltonian
operators are related to linear line
congruences, while third-order operators
and compatible first-order operators lead
to line complexes. Pl\"ucker coordinates
place these objects inside
$$
\mathbb{P}(\Lambda^2V).
$$
Compatibility of Hamiltonian operators is
thereby translated into compatibility of
projective varieties.

The broad point is that many familiar
integrable systems carry algebraic
varieties that organize their Hamiltonian
structures. Conversely, suitable
projective data can be used to construct
integrable systems. The Kaup--Broer system
was mentioned as an example arising from
this mechanism.

\section{PDE Without Tears for Geometry
Students}
\label{section-pde-without-tears}

\medskip\noindent
{\bf Speaker.}

Raffaele Vitolo,
Universit\`a del Salento.

\medskip\noindent
{\bf Seminar information.}

Semin\'ario das Sextas, PUC-Rio,
November 29, 2024.

The geometric study of differential
equations goes back to Lie and was put
into its modern form by Ehresmann. Its
basic idea is to regard derivatives as
coordinates on a larger space.

Let $N$ have dimension $n$ and let $M$
have dimension $m$. A field is a map
$$
u:U\longrightarrow M,
\qquad U\subset N.
$$
Locally it has independent coordinates
$x^1,\ldots,x^n$ and dependent
coordinates $u^1,\ldots,u^m$. It can
equivalently be regarded as the section
$$
s:N\longrightarrow N\times M,
\qquad s(x)=(x,u(x)).
$$

\medskip\noindent
{\bf Contact and jet spaces.}

\begin{definition}
\label{definition-contact-order}
Two local sections $s_1$ and $s_2$ have
\textit{contact of order $r$} at $x$ if
their derivatives at $x$ agree through
order $r$.
\end{definition}

Thus contact of order zero means that the
sections have the same value. Contact of
order one means in addition that their
graphs have the same tangent space.

More generally, let $E$ be a manifold of
dimension $n+m$, and let $L_1,L_2\subset E$
be submanifolds of dimension $n$ through
$z$. After expressing both as graphs in a
common chart, one can compare their
derivatives. The resulting notion of
contact is independent of the chart.

\begin{definition}
\label{definition-jet-space}
The \textit{$r$-jet} of $L$ at $z$ is its
equivalence class under contact of order
$r$, denoted by $[L]^r_z$. The
\textit{jet space} is
$$
J^r(E,n)=
\bigcup_{z\in E}
\left\{[L]^r_z\right\}.
$$
\end{definition}

There are natural projections
$$
\cdots\longrightarrow J^2(E,n)
\longrightarrow J^1(E,n)
\longrightarrow E.
$$
The first jet space records tangent
$n$-planes and is therefore a
Grassmannian bundle. Higher jet spaces
record higher-order Taylor data.

Suppose now that $E\to N$ is a bundle. A
local section has coordinates
$$
(x^\lambda,u^i),
$$
and its $r$-jet has coordinates
$$
(x^\lambda,u^i,u^i_\sigma),
\qquad |\sigma|\leq r.
$$
Here $\sigma$ is a multi-index and
$u^i_\sigma$ represents the corresponding
partial derivative of $u^i$.

\begin{definition}
\label{definition-pde-jet-space}
A \textit{partial differential equation of
order $r$} is a submanifold
$$
\mathcal{E}\subset J^r(E,n).
$$
\end{definition}

A system given locally by equations
$$
F^a(x^\lambda,u^i,u^i_\sigma)=0
$$
defines such a submanifold whenever the
usual regularity conditions hold.

\begin{definition}
\label{definition-prolongation-section}
The \textit{$r$-th prolongation} of a
section $s$ is
$$
j^rs:N\longrightarrow J^r(E,n),
\qquad x\longmapsto[s(N)]^r_{s(x)}.
$$
The section is a \textit{solution} of
$\mathcal{E}$ if
$$
j^rs(N)\subset\mathcal{E}.
$$
\end{definition}

\begin{example}
\label{example-first-derivative-one}
For $n=m=1$, the equation
$$
u_x=1
$$
is a surface in $J^1(\mathbb{R}^2,1)$.
Its solutions are
$$
u(x)=x+c.
$$
Their first prolongations lie in that
surface.
\end{example}

\medskip\noindent
{\bf The inviscid Burgers equation.}

\begin{example}
\label{example-inviscid-burgers}
The Riemann, Hopf, or inviscid Burgers
equation is
$$
u_t+uu_x=0.
$$
Along a characteristic curve one has
$$
\frac{d}{dt}u(t,x(t))=0,
\qquad
\frac{dx}{dt}=u.
$$
Consequently, if the initial condition is
$$
u(0,x)=f(x),
$$
then the solution is determined implicitly
by
$$
u(t,x)=f(x-tu(t,x)).
$$

Set
$$
\xi=x-tu,
\qquad
v(t,\xi)=u(t,x).
$$
Away from points where the coordinate
change is singular, Burgers' equation
becomes
$$
v_t=0.
$$
The singularity occurs when
$$
1+tf'(\xi)=0.
$$
Thus characteristics can intersect and a
smooth solution can develop a shock in
finite time.
\end{example}

This example suggests a general strategy:
find symmetries and conservation laws, and
then seek a change of variables that
simplifies or linearizes the equation.

\medskip\noindent
{\bf The Cartan distribution.}

The images of prolonged sections determine
distinguished tangent directions on a jet
space. These directions form the geometric
structure that must be preserved by any
transformation of differential equations.

\begin{definition}
\label{definition-cartan-distribution}
The \textit{Cartan distribution} on
$J^r(E,n)$ is the distribution spanned by
tangent spaces to prolonged submanifolds,
together with the vertical directions of
the projection to $J^{r-1}(E,n)$.
\end{definition}

For a bundle $E\to N$, introduce the total
derivatives
$$
D_\lambda=
\frac{\partial}{\partial x^\lambda}
+
\sum_{i,|\sigma|<r}
u^i_{\sigma+\lambda}
\frac{\partial}{\partial u^i_\sigma}.
$$
The Cartan distribution is locally spanned
by the $D_\lambda$ and by
$$
\frac{\partial}{\partial u^i_\tau},
\qquad |\tau|=r.
$$
The total derivatives commute:
$$
[D_\lambda,D_\mu]=0.
$$
They differentiate functions on jet space
as if all jet coordinates arose from an
actual section.

\medskip\noindent
{\bf Contact transformations and
symmetries.}

\begin{definition}
\label{definition-contact-transformation}
A \textit{contact transformation} of a jet
space is a local diffeomorphism preserving
the Cartan distribution.
\end{definition}

When there is more than one dependent
variable, the Lie--B\"acklund theorem says
that every contact transformation of a
finite-order jet space is the prolongation
of a transformation of $E$. With one
dependent variable, genuine contact
transformations already occur on
$J^1(E,n)$.

Contact transformations can sometimes
linearize nonlinear equations. A model
case is the Monge--Amp\`ere equation
$$
u_{xx}u_{yy}-u_{xy}^2=1,
$$
which may be studied using a Legendre
transformation.

\begin{definition}
\label{definition-symmetry-pde}
A \textit{symmetry} of a differential
equation $\mathcal{E}$ is a contact
transformation preserving $\mathcal{E}$.
An infinitesimal symmetry is a vector field
tangent to $\mathcal{E}$ whose flow
preserves the restricted Cartan
distribution.
\end{definition}

The bracket of two infinitesimal symmetries
is again an infinitesimal symmetry. For an
ordinary differential equation, the
Lie--Bianchi theorem says that a sufficiently
large solvable Lie algebra of symmetries
allows the equation to be solved by
quadratures. For PDEs, symmetries can reduce
the number of variables or turn a PDE into
an ODE.

\medskip\noindent
{\bf Prolongations and linearization.}

Let
$$
X=\xi^\lambda
\frac{\partial}{\partial x^\lambda}
+\phi^i\frac{\partial}{\partial u^i}
$$
be a vector field on $E$. Its characteristic
is
$$
Q^i=\phi^i-u^i_\lambda\xi^\lambda.
$$
Its prolongation to the infinite jet space
can be written as
$$
X^{(\infty)}=
\xi^\lambda D_\lambda
+
\sum_{i,\sigma}D_\sigma(Q^i)
\frac{\partial}{\partial u^i_\sigma}.
$$
This separates the motion along the base
from the evolutionary, or vertical, part.

Suppose the equation is defined by
$$
F^a(x^\lambda,u^i,u^i_\sigma)=0.
$$
Its linearization is the differential
operator
$$
\ell_F(Q)^a=
\sum_{i,\sigma}
\frac{\partial F^a}
{\partial u^i_\sigma}D_\sigma(Q^i).
$$
A generalized symmetry has characteristic
$Q$ satisfying
$$
\ell_F(Q)=0
$$
on the equation and all its differential
consequences.

\begin{example}
\label{example-kdv-linearization}
For the Korteweg--de Vries equation
$$
F=u_t+uu_x+u_{xxx}=0,
$$
the linearization is
$$
\ell_F(Q)=
D_tQ+uD_xQ+u_xQ+D_x^3Q.
$$
Time translation is a symmetry, represented
by the characteristic $Q=u_t$. Verifying
this means substituting the equation and
its total derivatives into
$\ell_F(u_t)$.
\end{example}

\medskip\noindent
{\bf Conservation laws.}

\begin{definition}
\label{definition-conservation-law}
For an equation in variables $t$ and $x$,
a \textit{conservation law} is a pair of
differential functions $\rho$ and $J$ such
that
$$
D_t\rho+D_xJ=0
$$
on every solution. The function $\rho$ is
the \textit{density}, and $J$ is the
\textit{flux}.
\end{definition}

Under suitable boundary conditions,
integration gives
$$
\frac{d}{dt}\int_a^b\rho\,dx
=-J(t,b)+J(t,a).
$$
If the boundary flux vanishes, the integral
of the density is constant in time.

Geometrically, conservation laws are
horizontal differential forms that become
closed after restriction to the equation.
They are the PDE analogue of first
integrals in ordinary differential
equations and are closely tied to
symmetries and Hamiltonian structures.

\section{Regularity Phenomena in
Nonlocal Equations}
\label{section-regularity-nonlocal}

\medskip\noindent
{\bf Conference information.}

\href{%
https://sites.google.com/view/%
60thanniversary-dmat-puc-rio/Home
}{International Conference Celebrating
the 60th Anniversary of the Department
of Mathematics at PUC-Rio}

August 12--14, 2026,
RDC Auditorium, PUC-Rio,
Rio de Janeiro.

\medskip\noindent
{\bf Speaker.}

Makson Santos,
Universidade de Lisboa.

\medskip\noindent
{\bf Date and time.}

August 13, 2026, 10:30--11:30.

\medskip\noindent
{\bf Title.}

Regularity Phenomena in Nonlocal
Equations.

\medskip\noindent
{\bf Abstract.}

The talk studies regularity for solutions
of fully nonlinear nonlocal equations in
elliptic and degenerate settings. In the
elliptic case, the goal is optimal
regularity for equations whose right-hand
side lies in $L^p$, with the regularity
space determined by $p$. In the degenerate
case, one obtains at least one viscosity
solution in
$C^{1,\alpha}_{\mathrm{loc}}$ for some
$\alpha\in(0,1)$. When the order of the
operator is sufficiently close to $2$,
one also obtains H\"older estimates for
the gradient of every viscosity solution.

\medskip\noindent
{\bf Notes.} Once we have a (weak) solution
of a PDE, we study what differentiability 
properties does the function have only
because it solves the given PDE.


\bibliography{my}
\bibliographystyle{amsalpha}

\end{document}
